La aplicación ha sido desarrollada íntegramente utilizando Visual Studio Code\footnote{\url{https://code.visualstudio.com/}} como editor de código principal, uno de los entornos de desarrollo más extendidos en la comunidad de desarrollo web gracias a su ligereza, su extensibilidad y su profunda integración con las tecnologías empleadas en este proyecto. A diferencia de los entornos de desarrollo integrado (IDE) orientados a un único ecosistema, Visual Studio Code actúa como un editor polivalente que, mediante un sistema de extensiones, proporciona soporte nativo para Python, JavaScript, HTML, CSS, y los lenguajes de plantillas utilizados tanto en el \textit{frontend} como en el \textit{backend}.

Para el desarrollo del \textit{backend} en Django, se configuró un entorno virtual de Python gestionado con uv\footnote{\url{https://github.com/astral-sh/uv}}, una herramienta moderna de gestión de dependencias y entornos que sustituye a pip y venv con un rendimiento notablemente superior. El \textit{frontend} con Vue 3 se desarrolló en el mismo espacio de trabajo, aprovechando la integración directa con el servidor de desarrollo de Vite, que proporciona recarga en caliente (\textit{HMR}) instantánea durante el proceso de creación de la interfaz. Por otro lado, la infraestructura del sistema se apoya en Docker mediante el despliegue de dos contenedores independientes, donde el primero alberga el motor de ejecución Piston para asegurar el aislamiento del código y el segundo gestiona la base de datos PostgreSQL para la persistencia de la información.

Todo el código fuente de la aplicación se encuentra disponible públicamente en un repositorio de GitHub (\url{https://github.com/ignaciogh33/SocratiCode}), lo que facilita su revisión, reproducción y posibles contribuciones futuras.
